\documentclass[aspectratio=169, 11pt]{beamer}
% ============================================================================
% Preambulo compartido para presentaciones Beamer
% Estadistica 2026-II - Pontificia Universidad Javeriana
% Incluir con: % ============================================================================
% Preambulo compartido para presentaciones Beamer
% Estadistica 2026-II - Pontificia Universidad Javeriana
% Incluir con: % ============================================================================
% Preambulo compartido para presentaciones Beamer
% Estadistica 2026-II - Pontificia Universidad Javeriana
% Incluir con: \input{../preambulo_beamer.tex} despues de \documentclass
% ============================================================================

\usepackage[T1]{fontenc}
\usepackage[utf8]{inputenc}
\usepackage[spanish]{babel}
\usepackage{amsmath, amssymb, amsfonts}
\usepackage{booktabs}
\usepackage{graphicx}
\usepackage{tikz}
\usetikzlibrary{shapes.geometric, positioning, arrows.meta, calc}
\usepackage{pgfplots}
\pgfplotsset{compat=1.18}
\usepackage{listings}
\usepackage{xcolor}
\usepackage{hyperref}
\usepackage{multicol}

% --- Tema Beamer (minimalista) ---
\usetheme{default}
\usecolortheme{default}
\usefonttheme{default}
\setbeamertemplate{navigation symbols}{}
\setbeamertemplate{caption}[numbered]
\setbeamertemplate{footline}{%
  \hspace{1em}{\tiny\url{https://github.com/polanco-jaime/Curso-Estadistica-2026-3}}\hfill\usebeamerfont{footline}\insertframenumber/\inserttotalframenumber\hspace{1em}\vspace{0.5em}%
}
\setbeamertemplate{itemize items}[circle]
\setbeamertemplate{enumerate items}[default]

% --- Colores institucionales (sutiles) ---
\definecolor{javeriana}{RGB}{0, 62, 105}
\definecolor{javerianalight}{RGB}{0, 105, 170}
\definecolor{codigobg}{RGB}{245, 245, 245}
\definecolor{codigofg}{RGB}{40, 40, 40}
\definecolor{comentario}{RGB}{100, 100, 100}
\definecolor{keyword}{RGB}{0, 100, 150}
\definecolor{string}{RGB}{180, 60, 30}

\setbeamercolor{title}{fg=javeriana}
\setbeamercolor{frametitle}{fg=javeriana}
\setbeamercolor{structure}{fg=javeriana}
\setbeamercolor{block title}{fg=white, bg=javeriana!75}
\setbeamercolor{block body}{bg=javeriana!5}
\setbeamercolor{block title alerted}{fg=white, bg=red!70!black}
\setbeamercolor{block body alerted}{bg=red!5}
\setbeamercolor{block title example}{fg=white, bg=green!40!black}
\setbeamercolor{block body example}{bg=green!5}
\setbeamercolor{item}{fg=javeriana}

\setbeamerfont{title}{series=\bfseries, size=\Large}
\setbeamerfont{frametitle}{series=\bfseries}

% --- Configuracion de listings para R ---
\lstset{
  language=R,
  basicstyle=\ttfamily\footnotesize,
  backgroundcolor=\color{codigobg},
  commentstyle=\color{comentario}\itshape,
  keywordstyle=\color{keyword}\bfseries,
  stringstyle=\color{string},
  numbers=none,
  frame=single,
  rulecolor=\color{javeriana!30},
  breaklines=true,
  showstringspaces=false,
  columns=flexible,
  morekeywords={library, ggplot, aes, geom_histogram, geom_boxplot,
    geom_point, geom_smooth, geom_density, geom_bar, geom_line,
    geom_vline, geom_hline, geom_errorbar, labs, theme_minimal,
    facet_wrap, scale_fill_manual, scale_color_manual,
    summarise, group_by, filter, mutate, select, arrange,
    read.csv, head, str, summary, mean, median, sd, var, cor,
    lm, t.test, chisq.test, wilcox.test, kruskal.test, aov,
    pnorm, qnorm, dnorm, rnorm, qt, pt, qchisq, pchisq, qf, pf,
    confint, coeftest, vcovHC, bptest, vif, cooks.distance,
    TukeyHSD, table, prop.table}
}

% --- Comandos personalizados ---
\newcommand{\R}{\texttt{R}}
\newcommand{\code}[1]{\texttt{\footnotesize #1}}
\newcommand{\variable}[1]{\texttt{#1}}
\newcommand{\paquete}[1]{\texttt{#1}}

% --- Informacion del curso ---
\institute{Pontificia Universidad Javeriana\\
  Facultad de Ciencias Econ\'omicas y Administrativas}
\date{2026-II}
 despues de \documentclass
% ============================================================================

\usepackage[T1]{fontenc}
\usepackage[utf8]{inputenc}
\usepackage[spanish]{babel}
\usepackage{amsmath, amssymb, amsfonts}
\usepackage{booktabs}
\usepackage{graphicx}
\usepackage{tikz}
\usetikzlibrary{shapes.geometric, positioning, arrows.meta, calc}
\usepackage{pgfplots}
\pgfplotsset{compat=1.18}
\usepackage{listings}
\usepackage{xcolor}
\usepackage{hyperref}
\usepackage{multicol}

% --- Tema Beamer (minimalista) ---
\usetheme{default}
\usecolortheme{default}
\usefonttheme{default}
\setbeamertemplate{navigation symbols}{}
\setbeamertemplate{caption}[numbered]
\setbeamertemplate{footline}{%
  \hspace{1em}{\tiny\url{https://github.com/polanco-jaime/Curso-Estadistica-2026-3}}\hfill\usebeamerfont{footline}\insertframenumber/\inserttotalframenumber\hspace{1em}\vspace{0.5em}%
}
\setbeamertemplate{itemize items}[circle]
\setbeamertemplate{enumerate items}[default]

% --- Colores institucionales (sutiles) ---
\definecolor{javeriana}{RGB}{0, 62, 105}
\definecolor{javerianalight}{RGB}{0, 105, 170}
\definecolor{codigobg}{RGB}{245, 245, 245}
\definecolor{codigofg}{RGB}{40, 40, 40}
\definecolor{comentario}{RGB}{100, 100, 100}
\definecolor{keyword}{RGB}{0, 100, 150}
\definecolor{string}{RGB}{180, 60, 30}

\setbeamercolor{title}{fg=javeriana}
\setbeamercolor{frametitle}{fg=javeriana}
\setbeamercolor{structure}{fg=javeriana}
\setbeamercolor{block title}{fg=white, bg=javeriana!75}
\setbeamercolor{block body}{bg=javeriana!5}
\setbeamercolor{block title alerted}{fg=white, bg=red!70!black}
\setbeamercolor{block body alerted}{bg=red!5}
\setbeamercolor{block title example}{fg=white, bg=green!40!black}
\setbeamercolor{block body example}{bg=green!5}
\setbeamercolor{item}{fg=javeriana}

\setbeamerfont{title}{series=\bfseries, size=\Large}
\setbeamerfont{frametitle}{series=\bfseries}

% --- Configuracion de listings para R ---
\lstset{
  language=R,
  basicstyle=\ttfamily\footnotesize,
  backgroundcolor=\color{codigobg},
  commentstyle=\color{comentario}\itshape,
  keywordstyle=\color{keyword}\bfseries,
  stringstyle=\color{string},
  numbers=none,
  frame=single,
  rulecolor=\color{javeriana!30},
  breaklines=true,
  showstringspaces=false,
  columns=flexible,
  morekeywords={library, ggplot, aes, geom_histogram, geom_boxplot,
    geom_point, geom_smooth, geom_density, geom_bar, geom_line,
    geom_vline, geom_hline, geom_errorbar, labs, theme_minimal,
    facet_wrap, scale_fill_manual, scale_color_manual,
    summarise, group_by, filter, mutate, select, arrange,
    read.csv, head, str, summary, mean, median, sd, var, cor,
    lm, t.test, chisq.test, wilcox.test, kruskal.test, aov,
    pnorm, qnorm, dnorm, rnorm, qt, pt, qchisq, pchisq, qf, pf,
    confint, coeftest, vcovHC, bptest, vif, cooks.distance,
    TukeyHSD, table, prop.table}
}

% --- Comandos personalizados ---
\newcommand{\R}{\texttt{R}}
\newcommand{\code}[1]{\texttt{\footnotesize #1}}
\newcommand{\variable}[1]{\texttt{#1}}
\newcommand{\paquete}[1]{\texttt{#1}}

% --- Informacion del curso ---
\institute{Pontificia Universidad Javeriana\\
  Facultad de Ciencias Econ\'omicas y Administrativas}
\date{2026-II}
 despues de \documentclass
% ============================================================================

\usepackage[T1]{fontenc}
\usepackage[utf8]{inputenc}
\usepackage[spanish]{babel}
\usepackage{amsmath, amssymb, amsfonts}
\usepackage{booktabs}
\usepackage{graphicx}
\usepackage{tikz}
\usetikzlibrary{shapes.geometric, positioning, arrows.meta, calc}
\usepackage{pgfplots}
\pgfplotsset{compat=1.18}
\usepackage{listings}
\usepackage{xcolor}
\usepackage{hyperref}
\usepackage{multicol}

% --- Tema Beamer (minimalista) ---
\usetheme{default}
\usecolortheme{default}
\usefonttheme{default}
\setbeamertemplate{navigation symbols}{}
\setbeamertemplate{caption}[numbered]
\setbeamertemplate{footline}{%
  \hspace{1em}{\tiny\url{https://github.com/polanco-jaime/Curso-Estadistica-2026-3}}\hfill\usebeamerfont{footline}\insertframenumber/\inserttotalframenumber\hspace{1em}\vspace{0.5em}%
}
\setbeamertemplate{itemize items}[circle]
\setbeamertemplate{enumerate items}[default]

% --- Colores institucionales (sutiles) ---
\definecolor{javeriana}{RGB}{0, 62, 105}
\definecolor{javerianalight}{RGB}{0, 105, 170}
\definecolor{codigobg}{RGB}{245, 245, 245}
\definecolor{codigofg}{RGB}{40, 40, 40}
\definecolor{comentario}{RGB}{100, 100, 100}
\definecolor{keyword}{RGB}{0, 100, 150}
\definecolor{string}{RGB}{180, 60, 30}

\setbeamercolor{title}{fg=javeriana}
\setbeamercolor{frametitle}{fg=javeriana}
\setbeamercolor{structure}{fg=javeriana}
\setbeamercolor{block title}{fg=white, bg=javeriana!75}
\setbeamercolor{block body}{bg=javeriana!5}
\setbeamercolor{block title alerted}{fg=white, bg=red!70!black}
\setbeamercolor{block body alerted}{bg=red!5}
\setbeamercolor{block title example}{fg=white, bg=green!40!black}
\setbeamercolor{block body example}{bg=green!5}
\setbeamercolor{item}{fg=javeriana}

\setbeamerfont{title}{series=\bfseries, size=\Large}
\setbeamerfont{frametitle}{series=\bfseries}

% --- Configuracion de listings para R ---
\lstset{
  language=R,
  basicstyle=\ttfamily\footnotesize,
  backgroundcolor=\color{codigobg},
  commentstyle=\color{comentario}\itshape,
  keywordstyle=\color{keyword}\bfseries,
  stringstyle=\color{string},
  numbers=none,
  frame=single,
  rulecolor=\color{javeriana!30},
  breaklines=true,
  showstringspaces=false,
  columns=flexible,
  morekeywords={library, ggplot, aes, geom_histogram, geom_boxplot,
    geom_point, geom_smooth, geom_density, geom_bar, geom_line,
    geom_vline, geom_hline, geom_errorbar, labs, theme_minimal,
    facet_wrap, scale_fill_manual, scale_color_manual,
    summarise, group_by, filter, mutate, select, arrange,
    read.csv, head, str, summary, mean, median, sd, var, cor,
    lm, t.test, chisq.test, wilcox.test, kruskal.test, aov,
    pnorm, qnorm, dnorm, rnorm, qt, pt, qchisq, pchisq, qf, pf,
    confint, coeftest, vcovHC, bptest, vif, cooks.distance,
    TukeyHSD, table, prop.table}
}

% --- Comandos personalizados ---
\newcommand{\R}{\texttt{R}}
\newcommand{\code}[1]{\texttt{\footnotesize #1}}
\newcommand{\variable}[1]{\texttt{#1}}
\newcommand{\paquete}[1]{\texttt{#1}}

% --- Informacion del curso ---
\institute{Pontificia Universidad Javeriana\\
  Facultad de Ciencias Econ\'omicas y Administrativas}
\date{2026-II}

\title{Sesión 3: Probabilidad, Distribución Normal y Muestreo}
\subtitle{De la probabilidad al Teorema Central del Límite}
\author{Prof.\ Jaime Polanco Jim\'enez}
\date{Vie 21 ago 2026}
\begin{document}

\begin{frame}
\titlepage
\end{frame}

\begin{frame}{Agenda}
\tableofcontents
\end{frame}

% ============================================================================
\section{Variables aleatorias}
% ============================================================================

\begin{frame}{Variable aleatoria}
\begin{block}{Definición}
Una \textbf{variable aleatoria} es una función que asigna un valor numérico a cada resultado de un experimento aleatorio.
\end{block}

\textbf{Notación}: usamos letras mayúsculas (X, Y, Z) para la variable aleatoria y minúsculas (x, y, z) para valores específicos.

\vspace{0.5em}
\begin{exampleblock}{Ejemplos}
\begin{itemize}
\item X = puntaje de inglés en Saber Pro (valores entre 0 y 200)
\item Y = número de estudiantes que aprueban un examen (valores 0, 1, 2, ..., n)
\item Z = tiempo de estudio en horas (valores $\geq 0$)
\end{itemize}
\end{exampleblock}

\vspace{0.5em}
\textbf{Tipos}:
\begin{itemize}
\item \textbf{Discreta}: toma valores separados, contables (ej: número de aprobados)
\item \textbf{Continua}: puede tomar cualquier valor en un intervalo (ej: puntaje, altura, tiempo)
\end{itemize}
\end{frame}

\begin{frame}{Función de probabilidad (variable discreta)}
\textbf{Variable discreta}: Función de probabilidad (pmf)
\[
P(X = x)
\]

Propiedades:
\begin{itemize}
\item $0 \leq P(X=x) \leq 1$
\item $\sum_{\text{todos } x} P(X=x) = 1$
\end{itemize}

\vspace{0.5em}
Ejemplo: X = resultado de un dado
\[
P(X=k) = \frac{1}{6}, \quad k=1,2,3,4,5,6
\]
\end{frame}

\begin{frame}{Función de densidad (variable continua)}
\textbf{Variable continua}: Función de densidad (pdf)
\[
f(x)
\]

Propiedades:
\begin{itemize}
\item $f(x) \geq 0$
\item $\int_{-\infty}^{\infty} f(x) dx = 1$
\item $P(a \leq X \leq b) = \int_a^b f(x) dx$
\end{itemize}

\vspace{0.5em}
\alert{Importante}: Para variables continuas,
\[
P(X = x) = 0
\]
Solo tiene sentido calcular $P(a \leq X \leq b)$
\end{frame}

\begin{frame}{Valor esperado y varianza}
\begin{block}{Valor esperado (media poblacional)}
El \textbf{valor esperado} de X, denotado $E(X)$ o $\mu$, es el promedio ponderado de todos los valores posibles.
\end{block}

\begin{columns}[T]
\column{0.48\textwidth}
\textbf{Discreta}:
\[
E(X) = \mu = \sum_{\text{todos } x} x \cdot P(X=x)
\]

\column{0.48\textwidth}
\textbf{Continua}:
\[
E(X) = \mu = \int_{-\infty}^{\infty} x \cdot f(x) dx
\]
\end{columns}

\vspace{1em}
\begin{block}{Varianza}
La \textbf{varianza} de X, denotada $\text{Var}(X)$ o $\sigma^2$, mide la dispersión alrededor de $\mu$.
\[
\text{Var}(X) = \sigma^2 = E[(X - \mu)^2] = E(X^2) - [E(X)]^2
\]
\textbf{Desviación estándar}: $\sigma = \sqrt{\sigma^2}$
\end{block}

\vspace{0.5em}
\alert{Nota}: $\mu$ y $\sigma$ son \textbf{parámetros poblacionales} (fijos, desconocidos). $\bar{x}$ y $s$ son \textbf{estadísticos muestrales} (variables, observables).
\end{frame}

\begin{frame}{Ejemplo: puntaje de inglés como variable aleatoria}
\begin{block}{Contexto}
Consideremos el puntaje de inglés en Saber Pro para estudiantes de NI como una variable aleatoria continua X.
\end{block}

\textbf{Interpretación probabilística}:
\begin{itemize}
\item X puede tomar valores entre 0 y 200
\item X tiene una distribución de probabilidad con media $\mu$ y desviación estándar $\sigma$
\item $P(X \leq 150)$ = proporción de estudiantes con puntaje $\leq$ 150
\item $P(160 \leq X \leq 180)$ = proporción de estudiantes entre 160 y 180
\end{itemize}

\vspace{0.5em}
\textbf{Relación con estadística descriptiva}:
\begin{itemize}
\item Si tenemos datos de \alert{toda la población} de estudiantes NI: $\mu$ y $\sigma$ se calculan directamente
\item Si tenemos una \alert{muestra}: estimamos $\mu$ con $\bar{x}$ y $\sigma$ con $s$
\end{itemize}

\vspace{0.5em}
\alert{Pregunta clave}: ¿Cuál es la distribución de X? ¿Es normal, sesgada, bimodal? Esto determina cómo calculamos probabilidades.
\end{frame}

% ============================================================================
\section{Distribución normal}
% ============================================================================

\begin{frame}{La distribución normal}
\begin{block}{Definición}
La \textbf{distribución normal} (o gaussiana) es la distribución de probabilidad continua más importante en estadística.
\end{block}

\textbf{Función de densidad}:
\[
f(x) = \frac{1}{\sigma\sqrt{2\pi}} e^{-\frac{1}{2}\left(\frac{x-\mu}{\sigma}\right)^2}, \quad -\infty < x < \infty
\]

\textbf{Notación}: $X \sim N(\mu, \sigma^2)$ (X sigue una distribución normal con media $\mu$ y varianza $\sigma^2$)

\vspace{0.5em}
\textbf{Parámetros}:
\begin{itemize}
\item $\mu$: media (centro de la distribución)
\item $\sigma$: desviación estándar (dispersión)
\item $\sigma^2$: varianza
\end{itemize}

\alert{Importante}: La normal está completamente determinada por $\mu$ y $\sigma$. Conociendo estos dos parámetros, podemos calcular cualquier probabilidad.
\end{frame}

\begin{frame}{Propiedades de la distribución normal}
\begin{enumerate}
\item \textbf{Forma de campana}: simétrica alrededor de $\mu$
\item \textbf{Unimodal}: un solo pico en $x = \mu$
\item \textbf{Media = Mediana = Moda} = $\mu$
\item \textbf{Asíntota}: las colas se acercan al eje x pero nunca lo tocan
\item \textbf{Área total} = 1 (integra a 1)
\item \textbf{Puntos de inflexión}: en $\mu \pm \sigma$
\item \textbf{Regla 68-95-99.7}:
  \begin{itemize}
  \item 68\% de los datos caen dentro de $\mu \pm 1\sigma$
  \item 95\% caen dentro de $\mu \pm 2\sigma$
  \item 99.7\% caen dentro de $\mu \pm 3\sigma$
  \end{itemize}
\end{enumerate}

\vspace{0.5em}
\begin{center}
\begin{tikzpicture}[scale=0.7]
\begin{axis}[
  domain=-4:4, samples=100, axis lines=middle,
  xlabel={$x$}, ylabel={$f(x)$}, ytick=\empty,
  xtick={-3,-2,-1,0,1,2,3},
  xticklabels={$\mu-3\sigma$,$\mu-2\sigma$,$\mu-\sigma$,$\mu$,$\mu+\sigma$,$\mu+2\sigma$,$\mu+3\sigma$},
  height=4cm, width=12cm, enlargelimits=false, clip=false
]
\addplot[very thick, javeriana] {exp(-x^2/2)/sqrt(2*pi)};
\end{axis}
\end{tikzpicture}
\end{center}
\end{frame}

\begin{frame}{Distribución normal estándar}
\begin{block}{Definición}
La \textbf{distribución normal estándar} es una distribución normal con $\mu = 0$ y $\sigma = 1$.
\[
Z \sim N(0, 1)
\]
\end{block}

\textbf{Función de densidad}:
\[
f(z) = \frac{1}{\sqrt{2\pi}} e^{-\frac{z^2}{2}}
\]

\textbf{Estandarización}: Cualquier variable normal $X \sim N(\mu, \sigma^2)$ se puede transformar a la normal estándar:
\[
Z = \frac{X - \mu}{\sigma} \sim N(0, 1)
\]

\vspace{0.5em}
\begin{alertblock}{¿Por qué es útil?}
\begin{itemize}
\item Todas las probabilidades se pueden calcular usando una sola tabla (tabla Z)
\item Permite comparar variables en diferentes escalas
\item Simplifica cálculos teóricos
\end{itemize}
\end{alertblock}
\end{frame}

\begin{frame}{Cálculo de probabilidades con la normal}
\textbf{Tipos de probabilidades}:

\begin{enumerate}
\item $P(X \leq a)$: probabilidad de que X sea menor o igual a $a$
  \begin{itemize}
  \item Área bajo la curva a la izquierda de $a$
  \end{itemize}
\item $P(X \geq a)$: probabilidad de que X sea mayor o igual a $a$
  \begin{itemize}
  \item $P(X \geq a) = 1 - P(X \leq a)$
  \end{itemize}
\item $P(a \leq X \leq b)$: probabilidad de que X esté entre $a$ y $b$
  \begin{itemize}
  \item $P(a \leq X \leq b) = P(X \leq b) - P(X \leq a)$
  \end{itemize}
\end{enumerate}

\vspace{0.5em}
\textbf{Procedimiento general}:
\begin{enumerate}
\item Identificar $\mu$ y $\sigma$ de la distribución
\item Estandarizar: $z = (x - \mu)/\sigma$
\item Buscar en tabla Z o usar software
\end{enumerate}
\end{frame}

\begin{frame}{Ejemplo: probabilidades con puntaje de inglés}
\begin{exampleblock}{Problema}
Supongamos que $X$ = puntaje de inglés en Saber Pro para NI sigue una distribución $N(155, 20^2)$ (media 155, desviación estándar 20).
\end{exampleblock}

\textbf{a) ¿Qué proporción de estudiantes obtiene puntaje menor o igual a 150?}

\[
P(X \leq 150) = P\left(Z \leq \frac{150 - 155}{20}\right) = P(Z \leq -0.25)
\]

Usando tabla Z o R: $P(Z \leq -0.25) = 0.4013$

\alert{Respuesta}: Aproximadamente 40\% de los estudiantes obtienen $\leq$ 150.

\vspace{0.5em}
\textbf{b) ¿Qué proporción supera 170?}

\[
P(X > 170) = 1 - P(X \leq 170) = 1 - P\left(Z \leq \frac{170-155}{20}\right) = 1 - P(Z \leq 0.75)
\]

$P(Z \leq 0.75) = 0.7734$, entonces $P(X > 170) = 1 - 0.7734 = 0.2266$

\alert{Respuesta}: Aproximadamente 23\% superan 170.
\end{frame}

\begin{frame}{Ejemplo: probabilidad en un intervalo}
\textbf{c) ¿Qué proporción obtiene entre 145 y 165?}

\[
P(145 \leq X \leq 165) = P(X \leq 165) - P(X \leq 145)
\]

Estandarizando:
\[
z_1 = \frac{145 - 155}{20} = -0.5, \quad z_2 = \frac{165 - 155}{20} = 0.5
\]

\[
P(-0.5 \leq Z \leq 0.5) = P(Z \leq 0.5) - P(Z \leq -0.5)
\]

\[
= 0.6915 - 0.3085 = 0.3830
\]

\alert{Respuesta}: Aproximadamente 38\% obtienen entre 145 y 165.

\vspace{0.5em}
\begin{center}
\begin{tikzpicture}[scale=0.6]
\begin{axis}[
  domain=-3:3, samples=100, axis lines=middle,
  xlabel={Puntaje}, ytick=\empty, xtick={-0.5,0,0.5},
  xticklabels={145, 155, 165}, height=4cm, width=9cm
]
\addplot[very thick, javeriana] {exp(-x^2/2)/sqrt(2*pi)};
\addplot[fill=javerianalight, opacity=0.5, domain=-0.5:0.5] {exp(-x^2/2)/sqrt(2*pi)} \closedcycle;
\end{axis}
\end{tikzpicture}
\end{center}
\end{frame}

\begin{frame}{Encontrar cuantiles (valores de x dada una probabilidad)}
\begin{block}{Problema inverso}
En lugar de calcular $P(X \leq x)$ dado $x$, ahora queremos encontrar $x$ dada la probabilidad.
\end{block}

\textbf{Ejemplo}: ¿Qué puntaje deja al 75\% de los estudiantes por debajo? (percentil 75)

Queremos encontrar $x$ tal que $P(X \leq x) = 0.75$

\vspace{0.5em}
\textbf{Procedimiento}:
\begin{enumerate}
\item Encontrar el valor $z$ en la normal estándar tal que $P(Z \leq z) = 0.75$
\item Usar tabla Z inversa o R: $z = 0.674$
\item Desestandardizar: $x = \mu + z\sigma = 155 + 0.674(20) = 168.48$
\end{enumerate}

\alert{Respuesta}: El percentil 75 es 168.48 puntos. El 75\% de estudiantes obtiene $\leq$ 168.48.

\vspace{0.5em}
\textbf{Otro ejemplo}: ¿Qué puntaje deja al 90\% por debajo?\\
$z = 1.282$, $x = 155 + 1.282(20) = 180.64$
\end{frame}

\begin{frame}[fragile]{Funciones de la normal en R (1/2)}
\begin{lstlisting}
# dnorm(): funcion de densidad f(x)
dnorm(155, mean = 155, sd = 20)
# [1] 0.01994711

# pnorm(): funcion de distribucion acumulada P(X <= x)
pnorm(150, mean = 155, sd = 20)  # P(X <= 150)
# [1] 0.4012937

pnorm(170, mean = 155, sd = 20)  # P(X <= 170)
# [1] 0.7733726

# P(X > 170)
1 - pnorm(170, mean = 155, sd = 20)
# [1] 0.2266274
\end{lstlisting}
\end{frame}

\begin{frame}[fragile]{Funciones de la normal en R (2/2)}
\begin{lstlisting}
# P(145 <= X <= 165)
pnorm(165, mean = 155, sd = 20) - pnorm(145, mean = 155, sd = 20)
# [1] 0.3829249

# qnorm(): cuantiles (inversa de pnorm)
qnorm(0.75, mean = 155, sd = 20)  # percentil 75
# [1] 168.4898

qnorm(0.90, mean = 155, sd = 20)  # percentil 90
# [1] 180.6311
\end{lstlisting}
\end{frame}

\begin{frame}[fragile]{rnorm(): generar valores aleatorios normales}
\begin{lstlisting}
# rnorm(): generar valores aleatorios de N(mu, sigma^2)
set.seed(123)  # para reproducibilidad
muestra <- rnorm(n = 1000, mean = 155, sd = 20)

# Verificar que la muestra tiene media ~ 155 y sd ~ 20
mean(muestra)
# [1] 154.9176
sd(muestra)
# [1] 19.82844

# Histograma con curva teorica superpuesta
library(ggplot2)
ggplot(data.frame(x = muestra), aes(x = x)) +
  geom_histogram(aes(y = after_stat(density)), bins = 30,
                 fill = "lightblue", color = "white") +
  stat_function(fun = dnorm, args = list(mean = 155, sd = 20),
                color = "red", linewidth = 1) +
  labs(title = "Muestra simulada vs distribucion teorica N(155, 20^2)",
       x = "Puntaje", y = "Densidad") +
  theme_minimal()
\end{lstlisting}
\end{frame}

\begin{frame}[fragile]{Comparar datos reales con curva normal (1/2)}
\begin{lstlisting}
# Cargar datos reales de ICFES
library(tidyverse)
icfes <- read_csv("datos/saber_pro_ni_2024.csv")

# Calcular media y sd muestrales
mu_hat <- mean(icfes$MOD_INGLES_PUNT, na.rm = TRUE)
sigma_hat <- sd(icfes$MOD_INGLES_PUNT, na.rm = TRUE)
\end{lstlisting}
\end{frame}

\begin{frame}[fragile]{Comparar datos reales con curva normal (2/2)}
\begin{lstlisting}
# Histograma de datos reales con curva normal teorica
ggplot(icfes, aes(x = MOD_INGLES_PUNT)) +
  geom_histogram(aes(y = after_stat(density)), bins = 40,
                 fill = "steelblue", color = "white", alpha = 0.7) +
  stat_function(fun = dnorm,
                args = list(mean = mu_hat, sd = sigma_hat),
                color = "red", linewidth = 1.2) +
  labs(title = "Distribucion real vs normal teorica",
       subtitle = paste("N(", round(mu_hat, 1), ",",
                        round(sigma_hat, 1), "^2)"),
       x = "Puntaje de ingles", y = "Densidad") +
  theme_minimal()

# ?Los datos se parecen a una distribucion normal?
shapiro.test(sample(icfes$MOD_INGLES_PUNT, 5000))
\end{lstlisting}
\end{frame}

% ============================================================================
\section{Distribuciones chi-cuadrado, t y F}
% ============================================================================

\begin{frame}{Distribución chi-cuadrado ($\chi^2$)}
\begin{block}{Definición}
Si $Z_1, Z_2, \ldots, Z_k$ son variables aleatorias independientes $N(0,1)$, entonces
\[
\chi^2 = Z_1^2 + Z_2^2 + \cdots + Z_k^2 \sim \chi^2(k)
\]
sigue una distribución chi-cuadrado con $k$ grados de libertad.
\end{block}

\textbf{Propiedades}:
\begin{itemize}
\item Sólo valores positivos ($\chi^2 \geq 0$)
\item Asimétrica (sesgada a la derecha)
\item Media: $E(\chi^2) = k$
\item Varianza: $\text{Var}(\chi^2) = 2k$
\item A medida que $k$ aumenta, se aproxima a la normal
\end{itemize}

\textbf{Uso principal}:
\begin{itemize}
\item Inferencia sobre varianzas: $(n-1)s^2/\sigma^2 \sim \chi^2(n-1)$
\item Pruebas de bondad de ajuste
\item Pruebas de independencia en tablas de contingencia
\end{itemize}
\end{frame}

\begin{frame}[fragile]{Distribución chi-cuadrado en R (1/2)}
\begin{lstlisting}
library(ggplot2)

# Graficar chi-cuadrado para diferentes grados de libertad
x <- seq(0, 20, by = 0.1)
df_chisq <- data.frame(
  x = rep(x, 4),
  y = c(dchisq(x, df = 2), dchisq(x, df = 5),
        dchisq(x, df = 10), dchisq(x, df = 15)),
  df = rep(c("df = 2", "df = 5", "df = 10", "df = 15"),
           each = length(x))
)
\end{lstlisting}
\end{frame}

\begin{frame}[fragile]{Distribución chi-cuadrado en R (2/2)}
\begin{lstlisting}
ggplot(df_chisq, aes(x = x, y = y, color = df)) +
  geom_line(linewidth = 1) +
  labs(title = "Distribucion chi-cuadrado para varios gl",
       x = "Valor", y = "Densidad", color = "Grados de libertad") +
  theme_minimal()

# Calcular probabilidades
pchisq(10, df = 5)  # P(chi^2 <= 10) con 5 gl
# [1] 0.9247726

# Cuantiles
qchisq(0.95, df = 5)  # percentil 95 con 5 gl
# [1] 11.0705
\end{lstlisting}
\end{frame}

\begin{frame}{Distribución t de Student}
\begin{block}{Definición}
Si $Z \sim N(0,1)$ y $V \sim \chi^2(k)$ son independientes, entonces
\[
t = \frac{Z}{\sqrt{V/k}} \sim t(k)
\]
sigue una distribución t de Student con $k$ grados de libertad.
\end{block}

\textbf{Propiedades}:
\begin{itemize}
\item Simétrica alrededor de 0 (como la normal)
\item Colas más pesadas que la normal (más probabilidad en valores extremos)
\item Media: $E(t) = 0$ (si $k > 1$)
\item Varianza: $\text{Var}(t) = k/(k-2)$ (si $k > 2$), mayor que la normal estándar
\item A medida que $k \to \infty$, $t(k) \to N(0,1)$
\item Para $k \geq 30$, es prácticamente idéntica a la normal
\end{itemize}

\textbf{Uso principal}:
\begin{itemize}
\item Inferencia sobre medias cuando $\sigma$ es desconocida: $\frac{\bar{X} - \mu}{s/\sqrt{n}} \sim t(n-1)$
\item Intervalos de confianza y pruebas de hipótesis para medias
\end{itemize}
\end{frame}

\begin{frame}{¿De dónde viene la distribución t?}
\textbf{Motivación histórica}:

Cuando queremos inferir sobre la media poblacional $\mu$:
\begin{itemize}
\item Si conocemos $\sigma$: $\frac{\bar{X} - \mu}{\sigma/\sqrt{n}} \sim N(0,1)$
\item Pero en la práctica, $\sigma$ es desconocido, usamos $s$
\item Reemplazar $\sigma$ por $s$ introduce incertidumbre adicional
\item William Gosset (seudónimo "Student") descubrió en 1908 que:
\end{itemize}

\[
\frac{\bar{X} - \mu}{s/\sqrt{n}} \sim t(n-1)
\]

\vspace{0.5em}
\textbf{Consecuencias}:
\begin{itemize}
\item Para muestras pequeñas ($n < 30$), la distribución t tiene colas más pesadas
\item Intervalos de confianza son más anchos (reflejan mayor incertidumbre)
\item A medida que $n$ crece, $s \to \sigma$ y $t \to N(0,1)$
\end{itemize}

\alert{Por eso usamos t en lugar de Z cuando estimamos $\sigma$ con $s$}
\end{frame}

\begin{frame}[fragile]{Distribución t en R (1/2)}
\begin{lstlisting}
# Comparar t con diferentes gl vs normal estandar
x <- seq(-4, 4, by = 0.05)
df_t <- data.frame(
  x = rep(x, 4),
  y = c(dt(x, df = 1), dt(x, df = 5), dt(x, df = 30), dnorm(x)),
  dist = rep(c("t(1)", "t(5)", "t(30)", "N(0,1)"),
             each = length(x))
)

ggplot(df_t, aes(x = x, y = y, color = dist, linetype = dist)) +
  geom_line(linewidth = 1) +
  scale_linetype_manual(values = c("solid", "solid", "solid", "dashed")) +
  labs(title = "Distribucion t vs Normal estandar",
       subtitle = "Nota: colas mas pesadas para t con gl bajos",
       x = "Valor", y = "Densidad") +
  theme_minimal()
\end{lstlisting}
\end{frame}

\begin{frame}[fragile]{Distribución t en R (2/2)}
\begin{lstlisting}
# Probabilidades y cuantiles
pt(2, df = 10)  # P(t <= 2) con 10 gl
# [1] 0.9634783

qt(0.975, df = 10)  # percentil 97.5% (para IC 95%)
# [1] 2.228139

qt(0.975, df = 30)  # con 30 gl
# [1] 2.042272

qnorm(0.975)  # normal estandar
# [1] 1.959964
\end{lstlisting}
\end{frame}

\begin{frame}{Distribución F de Fisher}
\begin{block}{Definición}
Si $U \sim \chi^2(k_1)$ y $V \sim \chi^2(k_2)$ son independientes, entonces
\[
F = \frac{U/k_1}{V/k_2} \sim F(k_1, k_2)
\]
sigue una distribución F con $k_1$ y $k_2$ grados de libertad (numerador y denominador).
\end{block}

\textbf{Propiedades}:
\begin{itemize}
\item Sólo valores positivos ($F \geq 0$)
\item Asimétrica (sesgada a la derecha)
\item Depende de dos parámetros: $k_1$ (gl numerador) y $k_2$ (gl denominador)
\end{itemize}

\textbf{Usos principales}:
\begin{itemize}
\item Comparar varianzas de dos poblaciones: $\frac{s_1^2/\sigma_1^2}{s_2^2/\sigma_2^2} \sim F(n_1-1, n_2-1)$
\item ANOVA (análisis de varianza): comparar medias de $k$ grupos
\item Pruebas de significancia global en regresión lineal múltiple
\end{itemize}
\end{frame}

\begin{frame}[fragile]{Distribución F en R (1/2)}
\begin{lstlisting}
# Graficar F con diferentes grados de libertad
x <- seq(0, 5, by = 0.05)
df_f <- data.frame(
  x = rep(x, 3),
  y = c(df(x, df1 = 5, df2 = 10),
        df(x, df1 = 10, df2 = 20),
        df(x, df1 = 20, df2 = 30)),
  df_pair = rep(c("F(5, 10)", "F(10, 20)", "F(20, 30)"),
                each = length(x))
)
\end{lstlisting}
\end{frame}

\begin{frame}[fragile]{Distribución F en R (2/2)}
\begin{lstlisting}
ggplot(df_f, aes(x = x, y = y, color = df_pair)) +
  geom_line(linewidth = 1) +
  labs(title = "Distribucion F para varios grados de libertad",
       x = "Valor", y = "Densidad",
       color = "Grados de libertad") +
  theme_minimal()

# Probabilidades
pf(2.5, df1 = 5, df2 = 20)  # P(F <= 2.5)
# [1] 0.9345867

# Cuantiles (valores criticos para pruebas)
qf(0.95, df1 = 5, df2 = 20)  # percentil 95
# [1] 2.710588
\end{lstlisting}
\end{frame}

% ============================================================================
\section{Muestreo y distribución muestral de $\bar{X}$}
% ============================================================================

\begin{frame}{Población vs muestra}
\begin{columns}[T]
\column{0.48\textwidth}
\textbf{Población}
\begin{itemize}
\item Conjunto completo de elementos de interés
\item Parámetros (fijos, usualmente desconocidos):
  \begin{itemize}
  \item Media: $\mu$
  \item Desviación estándar: $\sigma$
  \item Proporción: $p$
  \end{itemize}
\item Notación: $N$ = tamaño poblacional
\end{itemize}

\column{0.48\textwidth}
\textbf{Muestra}
\begin{itemize}
\item Subconjunto de la población
\item Estadísticos (variables, observables):
  \begin{itemize}
  \item Media: $\bar{x}$
  \item Desviación estándar: $s$
  \item Proporción: $\hat{p}$
  \end{itemize}
\item Notación: $n$ = tamaño muestral
\end{itemize}
\end{columns}

\vspace{1em}
\begin{block}{Objetivo de la estadística inferencial}
Usar estadísticos muestrales ($\bar{x}$, $s$, $\hat{p}$) para hacer inferencias sobre parámetros poblacionales ($\mu$, $\sigma$, $p$).
\end{block}

\vspace{0.5em}
\begin{exampleblock}{Ejemplo}
Población: todos los estudiantes de NI en Colombia ($N \approx 50,000$)\\
Muestra: 500 estudiantes seleccionados aleatoriamente de la base ICFES\\
Estimamos $\mu$ (puntaje medio poblacional) con $\bar{x}$ (puntaje medio muestral)
\end{exampleblock}
\end{frame}

\begin{frame}{Muestreo aleatorio simple}
\begin{block}{Definición}
En el \textbf{muestreo aleatorio simple} (MAS), cada elemento de la población tiene la misma probabilidad de ser seleccionado, y cada muestra del mismo tamaño tiene igual probabilidad.
\end{block}

\textbf{Supuestos clave}:
\begin{enumerate}
\item Las observaciones son \alert{independientes} entre sí
\item Cada observación proviene de la \alert{misma distribución} (población)
\item Notación: $X_1, X_2, \ldots, X_n$ son i.i.d. (independientes e idénticamente distribuidas)
\end{enumerate}

\vspace{0.5em}
\textbf{Consecuencias}:
\begin{itemize}
\item Evita sesgos de selección
\item Permite generalizar de la muestra a la población
\item Fundamento matemático de la inferencia estadística
\end{itemize}

\vspace{0.5em}
\alert{En la práctica}: muestreo aleatorio perfecto es difícil. A veces usamos muestreo estratificado, por conglomerados, etc., pero MAS es el punto de partida teórico.
\end{frame}

\begin{frame}{Distribución muestral de $\bar{X}$}
\begin{block}{Concepto fundamental}
La \textbf{distribución muestral de $\bar{X}$} es la distribución de probabilidad de todas las posibles medias muestrales de tamaño $n$ que podríamos obtener de una población.
\end{block}

\textbf{Experimento mental}:
\begin{enumerate}
\item Tomar una muestra aleatoria de tamaño $n$ y calcular $\bar{x}$
\item Devolver la muestra a la población
\item Repetir pasos 1-2 infinitas veces
\item La distribución de todos esos $\bar{x}$ es la distribución muestral de $\bar{X}$
\end{enumerate}

\vspace{0.5em}
\textbf{Propiedades} (para cualquier población con media $\mu$ y varianza $\sigma^2$):
\[
E(\bar{X}) = \mu
\]
\[
\text{Var}(\bar{X}) = \frac{\sigma^2}{n}, \quad \text{DE}(\bar{X}) = \frac{\sigma}{\sqrt{n}}
\]

\alert{Error estándar}: $\text{SE}(\bar{X}) = \sigma/\sqrt{n}$ (desviación estándar de $\bar{X}$)
\end{frame}

\begin{frame}{Interpretación del error estándar}
\[
\text{SE}(\bar{X}) = \frac{\sigma}{\sqrt{n}}
\]

\textbf{Mensajes clave}:
\begin{enumerate}
\item \alert{$\bar{X}$ es menos variable que $X$}
  \begin{itemize}
  \item $\text{DE}(X) = \sigma$, pero $\text{DE}(\bar{X}) = \sigma/\sqrt{n}$
  \item Promedios varían menos que valores individuales
  \end{itemize}
\item \alert{El error estándar disminuye con $\sqrt{n}$, no con $n$}
  \begin{itemize}
  \item Para reducir SE a la mitad, necesitamos 4 veces más observaciones
  \item Rendimientos decrecientes del tamaño muestral
  \end{itemize}
\item \alert{Más grande la muestra, más precisa la estimación de $\mu$}
  \begin{itemize}
  \item $n = 25$: $\text{SE} = \sigma/5$
  \item $n = 100$: $\text{SE} = \sigma/10$
  \item $n = 400$: $\text{SE} = \sigma/20$
  \end{itemize}
\end{enumerate}

\begin{exampleblock}{Ejemplo}
Si $\sigma = 20$ puntos, entonces con $n=100$: $\text{SE} = 20/10 = 2$ puntos.\\
La media muestral típicamente difiere de $\mu$ en aproximadamente 2 puntos.
\end{exampleblock}
\end{frame}

% ============================================================================
\section{Teorema Central del Límite}
% ============================================================================

\begin{frame}{Teorema Central del Límite (TCL)}
\begin{block}{Teorema Central del Límite}
Si $X_1, X_2, \ldots, X_n$ es una muestra aleatoria de tamaño $n$ de una población con media $\mu$ y desviación estándar $\sigma$ (finitos), entonces para $n$ suficientemente grande:
\[
\bar{X} \sim N\left(\mu, \frac{\sigma^2}{n}\right) \quad \text{aproximadamente}
\]
O equivalentemente:
\[
Z = \frac{\bar{X} - \mu}{\sigma/\sqrt{n}} \sim N(0, 1) \quad \text{aproximadamente}
\]
\end{block}

\vspace{0.5em}
\textbf{Mensajes clave}:
\begin{enumerate}
\item \alert{No importa la forma de la distribución original}: puede ser sesgada, bimodal, uniforme, etc.
\item \alert{Para $n$ grande, $\bar{X}$ es aproximadamente normal}
\item \alert{"Suficientemente grande"}: generalmente $n \geq 30$, pero depende de la forma original
  \begin{itemize}
  \item Si la población es normal: $\bar{X}$ es exactamente normal para cualquier $n$
  \item Si la población es simétrica: $n \approx 15$ suele ser suficiente
  \item Si la población es muy sesgada: puede necesitarse $n \geq 50$ o más
  \end{itemize}
\end{enumerate}
\end{frame}

\begin{frame}{¿Por qué el TCL es tan importante?}
\begin{alertblock}{El TCL es el teorema más importante de la estadística}
\end{alertblock}

\textbf{Razones}:
\begin{enumerate}
\item \textbf{Fundamento de la inferencia}:
  \begin{itemize}
  \item Intervalos de confianza para medias
  \item Pruebas de hipótesis sobre medias
  \item Tamaño de muestra para precisión deseada
  \end{itemize}
\item \textbf{Universalidad}:
  \begin{itemize}
  \item Aplica a cualquier población (con media y varianza finitas)
  \item No necesitamos conocer la distribución poblacional
  \end{itemize}
\item \textbf{Justifica el uso de la distribución normal}:
  \begin{itemize}
  \item Muchos fenómenos naturales y sociales son sumas o promedios
  \item Por el TCL, tienden a ser normales
  \end{itemize}
\item \textbf{Permite aproximaciones}:
  \begin{itemize}
  \item Binomial $\to$ Normal (para $n$ grande)
  \item Poisson $\to$ Normal (para $\lambda$ grande)
  \end{itemize}
\end{enumerate}

\alert{Sin el TCL, la inferencia estadística sería muchísimo más complicada}
\end{frame}

\begin{frame}[fragile]{Simulación TCL paso 1: población no normal (1/2)}
\begin{lstlisting}
library(tidyverse)

# Cargar datos Saber 11 (poblacion)
# Supongamos que tenemos puntajes de Saber 11 Ingles (no normales)
saber11 <- read_csv("datos/saber_11_2023.csv")

# Parametros poblacionales
mu <- mean(saber11$PUNT_INGLES, na.rm = TRUE)
sigma <- sd(saber11$PUNT_INGLES, na.rm = TRUE)
cat("Media poblacional:", round(mu, 2), "\n")
cat("DE poblacional:", round(sigma, 2), "\n")

# Verificar asimetria
library(moments)
skewness(saber11$PUNT_INGLES, na.rm = TRUE)
# [1] 0.42 (positiva, sesgada a la derecha)
\end{lstlisting}
\end{frame}

\begin{frame}[fragile]{Simulación TCL paso 1: población no normal (2/2)}
\begin{lstlisting}
# Graficar la distribucion poblacional
ggplot(saber11, aes(x = PUNT_INGLES)) +
  geom_histogram(aes(y = after_stat(density)), bins = 50,
                 fill = "steelblue", color = "white") +
  geom_density(color = "red", linewidth = 1) +
  labs(title = "Distribucion poblacional de Saber 11 Ingles",
       subtitle = "Nota: NO es normal (sesgada a la derecha)",
       x = "Puntaje", y = "Densidad") +
  theme_minimal()
\end{lstlisting}
\end{frame}

\begin{frame}[fragile]{Simulación TCL paso 2: muestrear (1/2)}
\begin{lstlisting}
# Funcion para tomar una muestra y calcular su media
tomar_muestra_y_promediar <- function(datos, n) {
  muestra <- sample(datos, size = n, replace = FALSE)
  mean(muestra, na.rm = TRUE)
}

# Simular 2000 muestras de tamano n = 30
set.seed(456)
n <- 30
num_simulaciones <- 2000

medias_n30 <- replicate(num_simulaciones,
                        tomar_muestra_y_promediar(saber11$PUNT_INGLES, n))
\end{lstlisting}
\end{frame}

\begin{frame}[fragile]{Simulación TCL paso 2: muestrear (2/2)}
\begin{lstlisting}
# Estadisticas de la distribucion muestral
mean(medias_n30)  # Debe estar cerca de mu
sd(medias_n30)    # Debe estar cerca de sigma/sqrt(n)

cat("Media de las medias muestrales:", round(mean(medias_n30), 2), "\n")
cat("DE de las medias muestrales:", round(sd(medias_n30), 2), "\n")
cat("SE teorico (sigma/sqrt(n)):", round(sigma/sqrt(n), 2), "\n")

# Comparacion
# mean(medias_n30) ~ mu  (insesgado)
# sd(medias_n30) ~ sigma/sqrt(n)  (error estandar)
\end{lstlisting}
\end{frame}

\begin{frame}[fragile]{Simulación TCL paso 3: visualizar normalidad (1/2)}
\begin{lstlisting}
# Graficar la distribucion de medias muestrales (n=30)
df_medias <- data.frame(medias = medias_n30)

ggplot(df_medias, aes(x = medias)) +
  geom_histogram(aes(y = after_stat(density)), bins = 40,
                 fill = "lightgreen", color = "white", alpha = 0.7) +
  stat_function(fun = dnorm,
                args = list(mean = mu, sd = sigma/sqrt(n)),
                color = "red", linewidth = 1.2) +
  labs(title = paste("Distribucion muestral de X-barra (n =", n, ")"),
       subtitle = "Curva roja: N(mu, sigma^2/n) teorica",
       x = "Media muestral", y = "Densidad") +
  theme_minimal()
\end{lstlisting}
\end{frame}

\begin{frame}[fragile]{Simulación TCL paso 3: visualizar normalidad (2/2)}
\begin{lstlisting}
# Observacion: !La distribucion de medias ES aproximadamente normal,
# aunque la poblacion original NO lo era!
# Esto es el Teorema Central del Limite en accion.

# Verificar normalidad con QQ-plot
ggplot(df_medias, aes(sample = medias)) +
  stat_qq() + stat_qq_line(color = "red") +
  labs(title = "QQ-plot de medias muestrales") +
  theme_minimal()
\end{lstlisting}
\end{frame}

\begin{frame}[fragile]{Simulación TCL paso 4: efecto de n (1/2)}
\begin{lstlisting}
# Repetir para diferentes tamanos de muestra: n = 10, 30, 100, 500
tamanos <- c(10, 30, 100, 500)
resultados <- list()

set.seed(789)
for (n in tamanos) {
  medias <- replicate(2000,
                      tomar_muestra_y_promediar(saber11$PUNT_INGLES, n))
  resultados[[paste0("n_", n)]] <- data.frame(
    medias = medias,
    n = paste("n =", n),
    se_teorico = sigma / sqrt(n)
  )
}

df_todas <- bind_rows(resultados)
\end{lstlisting}
\end{frame}

\begin{frame}[fragile]{Simulación TCL paso 4: efecto de n (2/2)}
\begin{lstlisting}
# Graficar todas juntas
ggplot(df_todas, aes(x = medias)) +
  geom_histogram(aes(y = after_stat(density)), bins = 30,
                 fill = "steelblue", alpha = 0.6) +
  facet_wrap(~ n, scales = "free", ncol = 2) +
  labs(title = "Efecto del tamano de muestra en la distribucion de X-barra",
       x = "Media muestral", y = "Densidad") +
  theme_minimal()

# Observacion: a mayor n, la distribucion es mas concentrada (menor SE)
# y mas normal
\end{lstlisting}
\end{frame}

\begin{frame}[fragile]{Comparación SE teórico vs simulado (1/2)}
\begin{lstlisting}
# Comparar SE teorico vs simulado para cada n
comparacion_se <- df_todas %>%
  group_by(n) %>%
  summarise(
    SE_simulado = sd(medias),
    SE_teorico = first(se_teorico)
  ) %>%
  mutate(diferencia = abs(SE_simulado - SE_teorico))

comparacion_se

# Resultado esperado:
#   n         SE_simulado  SE_teorico  diferencia
#   n = 10         3.15        3.16        0.01
#   n = 30         1.83        1.83        0.00
#   n = 100        1.00        1.00        0.00
#   n = 500        0.45        0.45        0.00
\end{lstlisting}
\end{frame}

\begin{frame}[fragile]{Comparación SE teórico vs simulado (2/2)}
\begin{lstlisting}
# Conclusion: el SE teorico (sigma/sqrt(n)) predice muy bien
# la variabilidad de X-barra

# Grafico de barras comparando SE
ggplot(comparacion_se, aes(x = n)) +
  geom_col(aes(y = SE_teorico), fill = "red", alpha = 0.5, width = 0.4) +
  geom_col(aes(y = SE_simulado), fill = "blue", alpha = 0.5,
           width = 0.4, position = position_nudge(x = 0.3)) +
  labs(title = "SE teorico vs simulado",
       y = "Error estandar") +
  theme_minimal()
\end{lstlisting}
\end{frame}

% ============================================================================
\section{Propiedades de estimadores}
% ============================================================================

\begin{frame}{¿Qué es un estimador?}
\begin{block}{Definición}
Un \textbf{estimador} es una regla o fórmula para calcular una aproximación de un parámetro poblacional desconocido a partir de datos muestrales.
\end{block}

\textbf{Notación}:
\begin{itemize}
\item Parámetro (desconocido): $\theta$ (puede ser $\mu$, $\sigma^2$, $p$, etc.)
\item Estimador (función de los datos): $\hat{\theta}$ (ej: $\bar{X}$, $s^2$, $\hat{p}$)
\item Estimación (valor numérico): $\hat{\theta} = 155.3$ puntos
\end{itemize}

\vspace{0.5em}
\begin{exampleblock}{Ejemplos de estimadores}
\begin{itemize}
\item Estimador de $\mu$: $\hat{\mu} = \bar{X} = \frac{1}{n}\sum X_i$
\item Estimador de $\sigma^2$: $\hat{\sigma}^2 = s^2 = \frac{1}{n-1}\sum(X_i - \bar{X})^2$
\item Estimador de $p$ (proporción): $\hat{p} = \frac{X}{n}$ (donde $X$ = número de éxitos)
\end{itemize}
\end{exampleblock}

\vspace{0.5em}
\alert{Pregunta}: ¿Cómo sabemos si un estimador es "bueno"? $\to$ Propiedades deseables
\end{frame}

\begin{frame}{Propiedad 1: Insesgamiento}
\begin{block}{Definición}
Un estimador $\hat{\theta}$ es \textbf{insesgado} si su valor esperado es igual al parámetro:
\[
E(\hat{\theta}) = \theta
\]
En promedio (sobre todas las muestras posibles), el estimador acierta al valor verdadero.
\end{block}

\textbf{Sesgo}: $\text{Sesgo}(\hat{\theta}) = E(\hat{\theta}) - \theta$

\vspace{0.5em}
\begin{exampleblock}{Ejemplos}
\begin{itemize}
\item $\bar{X}$ es insesgado para $\mu$: $E(\bar{X}) = \mu$
\item $s^2 = \frac{1}{n-1}\sum(X_i-\bar{X})^2$ es insesgado para $\sigma^2$
\item $\hat{\sigma}^2 = \frac{1}{n}\sum(X_i-\bar{X})^2$ es \alert{sesgado} (subestima $\sigma^2$)
  \begin{itemize}
  \item $E(\hat{\sigma}^2) = \frac{n-1}{n}\sigma^2 < \sigma^2$
  \item Por eso usamos $n-1$ en el denominador
  \end{itemize}
\end{itemize}
\end{exampleblock}

\vspace{0.5em}
\alert{Insesgamiento no garantiza precisión en una muestra individual, pero sí "en promedio"}
\end{frame}

\begin{frame}{Propiedad 2: Eficiencia}
\begin{block}{Definición}
Entre dos estimadores insesgados de $\theta$, el más \textbf{eficiente} es el que tiene menor varianza.
\end{block}

\[
\text{Si } E(\hat{\theta}_1) = E(\hat{\theta}_2) = \theta \text{ y } \text{Var}(\hat{\theta}_1) < \text{Var}(\hat{\theta}_2), \text{ entonces } \hat{\theta}_1 \text{ es más eficiente}
\]

\vspace{0.5em}
\textbf{Interpretación}:
\begin{itemize}
\item Menor varianza = estimaciones más consistentes entre muestras
\item Menor varianza = intervalos de confianza más estrechos
\end{itemize}

\vspace{0.5em}
\begin{exampleblock}{Ejemplo}
Para estimar $\mu$ de una distribución normal:
\begin{itemize}
\item Media muestral: $\text{Var}(\bar{X}) = \sigma^2/n$
\item Mediana muestral: $\text{Var}(\text{mediana}) \approx 1.57 \sigma^2/n$
\end{itemize}
$\bar{X}$ es más eficiente que la mediana (para normales).
\end{exampleblock}
\end{frame}

\begin{frame}{Propiedad 3: Consistencia}
\begin{block}{Definición}
Un estimador $\hat{\theta}_n$ es \textbf{consistente} si converge al parámetro verdadero cuando el tamaño de muestra tiende a infinito:
\[
\hat{\theta}_n \xrightarrow{P} \theta \quad \text{cuando } n \to \infty
\]
\end{block}

\textbf{Interpretación}:
\begin{itemize}
\item Con más datos, el estimador se acerca más y más al valor verdadero
\item Probabilidad de estar cerca de $\theta$ tiende a 1
\end{itemize}

\vspace{0.5em}
\textbf{Condición suficiente para consistencia}:
\[
\lim_{n \to \infty} E(\hat{\theta}_n) = \theta \quad \text{y} \quad \lim_{n \to \infty} \text{Var}(\hat{\theta}_n) = 0
\]

\vspace{0.5em}
\begin{exampleblock}{Ejemplos}
\begin{itemize}
\item $\bar{X}$ es consistente: $E(\bar{X}) = \mu$ y $\text{Var}(\bar{X}) = \sigma^2/n \to 0$
\item $s^2$ es consistente para $\sigma^2$
\end{itemize}
\end{exampleblock}

\alert{Todos los estimadores comunes que usamos son consistentes}
\end{frame}

\begin{frame}{Resumen: propiedades de estimadores}
\begin{table}
\centering
\small
\begin{tabular}{lp{8cm}}
\toprule
\textbf{Propiedad} & \textbf{Significado} \\
\midrule
Insesgamiento & En promedio, el estimador acierta ($E(\hat{\theta}) = \theta$) \\
\addlinespace
Eficiencia & Menor varianza entre estimadores insesgados (más preciso) \\
\addlinespace
Consistencia & Converge al parámetro cuando $n \to \infty$ \\
\bottomrule
\end{tabular}
\end{table}

\vspace{1em}
\textbf{Estimadores comunes y sus propiedades}:
\begin{itemize}
\item $\bar{X}$ para $\mu$: insesgado, eficiente (para normales), consistente
\item $s^2$ para $\sigma^2$: insesgado, consistente
\item $\hat{p} = X/n$ para $p$: insesgado, consistente
\end{itemize}

\vspace{0.5em}
\begin{alertblock}{Trade-off sesgo-varianza}
A veces un estimador ligeramente sesgado puede tener mucho menor varianza, y ser preferible en la práctica (ej: regresión ridge). Pero en este curso usamos estimadores insesgados.
\end{alertblock}
\end{frame}

% ============================================================================
\section{Resumen y conexión}
% ============================================================================

\begin{frame}{Resumen de la sesión}
\textbf{Conceptos clave cubiertos}:

\begin{enumerate}
\item \textbf{Variables aleatorias}: discretas vs continuas, E(X), Var(X)
\item \textbf{Distribución normal}: $N(\mu, \sigma^2)$, propiedades, cálculo de probabilidades
  \begin{itemize}
  \item \code{pnorm()}, \code{qnorm()}, \code{dnorm()}, \code{rnorm()}
  \end{itemize}
\item \textbf{Distribuciones t, chi-cuadrado, F}: de dónde vienen, para qué sirven
\item \textbf{Muestreo}: población vs muestra, distribución muestral de $\bar{X}$
  \begin{itemize}
  \item $E(\bar{X}) = \mu$, $\text{SE}(\bar{X}) = \sigma/\sqrt{n}$
  \end{itemize}
\item \textbf{Teorema Central del Límite}: $\bar{X} \sim N(\mu, \sigma^2/n)$ para $n$ grande
  \begin{itemize}
  \item Fundamento de toda la inferencia estadística
  \end{itemize}
\item \textbf{Propiedades de estimadores}: insesgamiento, eficiencia, consistencia
\end{enumerate}
\end{frame}

\begin{frame}{La distribución normal + TCL = fundamento de la inferencia}
\begin{alertblock}{Conexión fundamental}
\end{alertblock}

\textbf{1. Muchas variables en la naturaleza son normales (o aproximadamente)}
\begin{itemize}
\item Alturas, pesos, puntajes de pruebas, errores de medición, etc.
\end{itemize}

\textbf{2. Incluso cuando la población NO es normal, $\bar{X}$ sí lo es (TCL)}
\begin{itemize}
\item Para $n \geq 30$, $\bar{X} \sim N(\mu, \sigma^2/n)$ aproximadamente
\end{itemize}

\textbf{3. Conociendo la distribución de $\bar{X}$, podemos}:
\begin{itemize}
\item Construir \alert{intervalos de confianza}: ¿entre qué valores está $\mu$ con 95\% de confianza?
\item Hacer \alert{pruebas de hipótesis}: ¿es $\mu = 150$ compatible con nuestros datos?
\item Calcular \alert{tamaños de muestra}: ¿cuántos estudiantes necesito para estimar $\mu$ con margen de error $\pm 2$ puntos?
\end{itemize}

\vspace{0.5em}
\alert{Sin el TCL, necesitaríamos conocer la distribución poblacional exacta para hacer inferencia. Gracias al TCL, no es necesario.}
\end{frame}

\begin{frame}{Próxima sesión: Intervalos de Confianza}
\textbf{Fecha}: próxima sesión (según calendario)

\textbf{Temas a cubrir}:
\begin{itemize}
\item Concepto de intervalo de confianza (IC)
\item IC para la media (varianza conocida y desconocida)
\item IC para la proporción
\item IC para la diferencia de medias (dos muestras)
\item Interpretación correcta de ICs
\item Determinación del tamaño de muestra
\end{itemize}

\vspace{1em}
\textbf{Preparación}:
\begin{itemize}
\item Repasar \code{qnorm()}, \code{qt()}, \code{pnorm()}, \code{pt()}
\item Asegurarse de entender el error estándar: $\text{SE} = \sigma/\sqrt{n}$
\item Repasar el TCL: $\bar{X} \sim N(\mu, \sigma^2/n)$ para $n$ grande
\end{itemize}

\vspace{0.5em}
\alert{El intervalo de confianza es la aplicación más directa del TCL y la distribución normal}
\end{frame}

\end{document}
